% =====================================================================
\section{Introduction}\label{sec:intro}
% =====================================================================

Contemporary video generation systems can synthesize visually compelling,
temporally coherent clips with synchronized audio. Perceptual plausibility,
however, is not evidence of physical validity: a failure may arise not only in
rendered appearance but in any observable consequence of a physical event.
Within the visual stream, entities may be omitted, trajectories may be
inconsistent with the forces implied by the scene, and interactions may violate
contact or causal constraints. Within the acoustic stream, an expected sound
may be absent, mistimed, or attributed to the wrong source; across modalities,
visual and acoustic observations may encode incompatible event timing or
causality. Thus an object may remain unsupported under gravity, a ball may pass
through a rigid wall, a visible collision may be silent, or its impact sound may
precede contact. These failures become especially consequential as generated
video moves beyond media synthesis into the training and evaluation loops of
embodied systems. Recent video world models have been used to synthesize
robot-training trajectories, support policy learning and evaluation, and reduce
the sim-to-real gap; stronger world models have also been associated with better
downstream policy
performance~\cite{jang2025dreamgen,huang2025enerverse,shang2025roboscape}.
Recently, evaluation has expanded beyond generic video quality to embodied task
correctness and rule-level physical
diagnosis~\cite{li2025worldmodelbench,deng2026rethinking,videophy2_2025,lin2026phyground},
while complementary work has introduced temporal localization and the evaluation
of acoustic or cross-modal physical consistency~\cite{physioneval2026,chinchure2026spotlight,xie2025phyavbench,cui2026avphys}.

\begin{figure}[t]
  \centering
  % =====================================================================
% fig_teaser.tex -- Figure 1 teaser: an opaque learned scalar judge versus
% PhyReAct's auditable, per-claim, timeline-grounded verdict.
%
% \input-able into main.tex, which loads adobe_research.cls + diffphy-support
% (hence the adobered / adobeink / adobegray colours and the sans family).
%
% REQUIRED IN THE MAIN PREAMBLE:
%     \usepackage{tikz}
%     \usetikzlibrary{arrows.meta,positioning,calc}
%
% Every style is declared inside the tikzpicture's own option list, so all
% definitions are local to this picture and cannot clash with the sibling
% schematics (fig_critic_pipeline.tex, fig_plan_graph.tex).
%
% CONTENT is illustrative of the method's output contract only -- the clip,
% the counts, and the windows are a worked schematic, not a measured run, in
% the same spirit as fig:plangraph. No accuracy or comparison number appears.
% =====================================================================
% The three TikZ libraries must be loaded in the PREAMBLE, not here: a figure
% environment forms a group, and TikZ defines a library's macros locally, so a
% \usetikzlibrary issued inside a figure does not survive to the next one.
\makeatletter
\@ifundefined{tikz@library@calc@loaded}{%
  \PackageError{fig_teaser}{%
    The preamble must contain\MessageBreak
    \string\usetikzlibrary{arrows.meta,positioning,calc}%
  }{Add that line after \string\usepackage{tikz}.}}{}
\makeatother

\begingroup
\providecommand{\PhyReAct}{PhyReAct}
\providecommand{\cfsub}[1]{{\scriptsize\color{adobegray}#1}}
\providecommand{\cfop}[1]{{\ttfamily\scriptsize\color{adobeink}#1}}

\begin{tikzpicture}[
    x=1cm, y=1cm,
    % --- shared box shape -------------------------------------------------
    cfbox/.style   = {rounded corners=2pt, line width=0.7pt, align=center,
                      inner sep=5pt, font=\sffamily\footnotesize,
                      text=adobeink},
    % --- LEARNED component: light gray fill --------------------------------
    cflearn/.style = {cfbox, draw=adobeink, fill=adobegray!12},
    % --- DETERMINISTIC component: white with an ink outline ----------------
    cfdet/.style   = {cfbox, draw=adobeink, fill=white},
    % --- what enters each evaluator ----------------------------------------
    cfdata/.style  = {cfbox, draw=adobegray, fill=white, inner sep=4pt,
                      dash pattern=on 2pt off 1.6pt, font=\sffamily\scriptsize},
    cfline/.style  = {draw=adobeink, line width=0.8pt},
    cfarr/.style   = {cfline, -{Straight Barb[length=4.6pt,width=5.2pt]}},
    cfnote/.style  = {font=\sffamily\scriptsize, text=adobegray,
                      inner sep=0pt, align=left},
    cfhead/.style  = {font=\sffamily\bfseries\footnotesize, text=adobegray,
                      inner sep=0pt, anchor=south west},
    % column-role captions over the three-stage grid
    cfcol/.style   = {font=\sffamily\scriptsize\bfseries, text=adobegray,
                      inner sep=0pt, anchor=south},
    % a filled dot where feed lines branch, so a bus reads as intentional
    cfdot/.style   = {circle, fill=adobeink, inner sep=0pt, minimum size=2.8pt},
    % --- claim-state badges (three-valued) ---------------------------------
    badge/.style   = {rounded corners=2pt, line width=0.7pt,
                      font=\sffamily\bfseries\scriptsize, inner sep=3pt,
                      align=center},
    bcon/.style    = {badge, draw=adobered, fill=adobered, text=adobeink},
    bsup/.style    = {badge, draw=adobeink, fill=adobeink, text=white},
    bunk/.style    = {badge, draw=adobegray, fill=white, text=adobegray},
  ]

  % =================================================================
  %  TOP BAND -- prior art: one learned scalar
  % =================================================================
  \node[cfhead] at (0,0.35) {PRIOR\ \ a learned scalar judge};

  \node[cfdata, text width=1.45cm, anchor=west]
    (in1) at (0,-0.85) {\textbf{Prompt}\\$+$ video};
  \node[cflearn, text width=3.5cm, anchor=west]
    (scorer) at (2.4,-0.85)
    {\textbf{Learned video--text scorer}\\[1pt]\cfsub{one forward pass}};
  \node[cfbox, draw=adobeink, fill=adobeink, text=white, minimum width=1.9cm,
        minimum height=0.95cm, font=\sffamily\bfseries\Large, anchor=west]
    (score) at (7.0,-0.85) {score};
  \node[cfnote, text width=5.2cm, anchor=west] (op1) at (9.1,-0.85)
    {one opaque score: it cannot say \emph{which} obligation failed,
     \emph{where} in the clip, or on \emph{what} evidence};

  \draw[cfarr] (in1.east)    -- (scorer.west);
  \draw[cfarr] (scorer.east) -- (score.west);

  % =================================================================
  %  DIVIDER
  % =================================================================
  \draw[cfline, dash pattern=on 2pt off 1.6pt, draw=adobegray]
    (0,-1.85) -- (14.1,-1.85);
  \node[fill=white, inner sep=2pt, text=adobegray,
        font=\sffamily\itshape\footnotesize] at (7.0,-1.85) {vs.};

  % =================================================================
  %  BOTTOM BAND -- ours: plan, then measure on the timeline
  % =================================================================
  \node[cfhead] at (0,-2.55)
    {\textbf{\PhyReAct{} (ours)}\ \ plan, then measure on the timeline};

  \node[cfdata, text width=1.45cm, anchor=west]
    (in2) at (0,-3.35) {\textbf{Prompt}\\$+$ video};
  \node[cflearn, text width=5.0cm, anchor=west]
    (planner) at (2.4,-3.35)
    {\textbf{Planner}\\[1pt]
     \cfsub{compiles the prompt into typed claims and an executable plan
     --- before any frame is read}};
  \draw[cfarr] (in2.east) -- (planner.west);

  % --- three claim rows on a fixed grid ----------------------------------
  % Uniform row pitch and a fixed minimum height on every box, so each box's
  % vertical centre sits exactly on its row baseline and every horizontal
  % arrow between columns is dead level. Column x-positions are shared.
  \def\rya{-5.15}\def\ryb{-6.30}\def\ryc{-7.45}
  \def\xC{0}\def\xS{4.7}\def\xB{9.5}
  \tikzset{gridbox/.style={minimum height=1.02cm, text depth=0pt}}

  % claim column (deterministic: white, ink outline)
  \node[cfdet, gridbox, text width=3.7cm, anchor=west] (c1) at (\xC,\rya)
    {\textbf{``two balls''}\\[1pt]\cfop{eq(count("ball"),\,2)}};
  \node[cfdet, gridbox, text width=3.7cm, anchor=west] (c2) at (\xC,\ryb)
    {\textbf{``bounces off the floor''}\\[1pt]\cfop{traj("ball",\,bounce)}};
  \node[cfdet, gridbox, text width=3.7cm, anchor=west] (c3) at (\xC,\ryc)
    {\textbf{``GO'' on the scoreboard}\\[1pt]\cfop{text("GO",\,scoreboard)}};

  % executor column (one frozen tool per row)
  \node[cflearn, gridbox, text width=3.4cm, anchor=west] (s1) at (\xS,\rya)
    {\textbf{Counting instrument}\\[1pt]\cfsub{returns count $=3$}};
  \node[cflearn, gridbox, text width=3.4cm, anchor=west] (s2) at (\xS,\ryb)
    {\textbf{Point tracker}\\[1pt]\cfsub{rebound after contact}};
  \node[cflearn, gridbox, text width=3.4cm, anchor=west] (s3) at (\xS,\ryc)
    {\textbf{OCR on localized frames}\\[1pt]\cfsub{glyphs illegible}};

  % three-valued claim state
  \node[bcon, gridbox, text width=2.9cm, anchor=west] (b1) at (\xB,\rya)
    {$\times$\, contradicted\\[1pt]\mdseries\scriptsize $1.2$--$2.4$\,s};
  \node[bsup, gridbox, text width=2.9cm, anchor=west] (b2) at (\xB,\ryb)
    {$\checkmark$\, supported};
  \node[bunk, gridbox, text width=2.9cm, anchor=west] (b3) at (\xB,\ryc)
    {\textbf{?}\, unknown\\[1pt]\mdseries\scriptsize $\to$ abstain};

  % column-role captions: name each stage (the middle column IS the executor)
  \node[cfcol] at ([yshift=3pt]c1.north) {claim $+$ typed check};
  \node[cfcol] at ([yshift=3pt]s1.north) {executor (one tool call)};
  \node[cfcol] at ([yshift=3pt]b1.north) {claim state};

  % per-row arrows: claim -> executor -> state, all level by construction
  \foreach \c/\s/\b in {c1/s1/b1, c2/s2/b2, c3/s3/b3}{
    \draw[cfarr] (\c.east) -- (\s.west);
    \draw[cfarr] (\s.east) -- (\b.west);
  }

  % --- planner fans out to the three claims -------------------------------
  % One stub drops from the planner into a manifold that runs in the clear gap
  % above the column captions, down a left-margin bus, and taps each claim on
  % its own row. The manifold sits ABOVE the caption row, so it never crosses
  % a claim box. Dots mark the taps.
  \def\xbus{-0.9}
  \coordinate (mani) at (\xbus,-4.15);            % manifold height: clear gap
  \draw[cfline] (planner.south) |- (mani);        % drop, then left onto the bus
  \draw[cfline] (mani) -- (\xbus,\ryc);           % vertical bus down the margin
  \foreach \y in {\rya,\ryb,\ryc}{
    \node[cfdot] at (\xbus,\y) {};
    \draw[cfarr] (\xbus,\y) -- (\xC,\y);          % level tap into each claim
  }

  % =================================================================
  %  DETERMINISTIC ROLL-UP + localized outputs
  % =================================================================
  \def\ryd{-9.05}
  \node[cfdet, text width=3.0cm, anchor=west] (roll) at (\xC,\ryd)
    {\textbf{Deterministic roll-up}\\[1pt]
     \cfsub{any one contradiction $\Rightarrow$ clip implausible}};
  \node[bcon, text width=2.6cm, anchor=west, font=\sffamily\bfseries\small,
        minimum height=0.9cm] (clip) at (4.4,\ryd)
    {clip:\\ IMPLAUSIBLE};
  \node[cfdet, text width=5.6cm, anchor=west] (pkt) at (8.4,\ryd)
    {\textbf{Localized evidence packet} $\mathcal{F}$\\[1pt]
     \cfsub{what failed $\cdot$ the supporting evidence $\cdot$
     where in the clip}};

  \draw[cfarr] (roll.east) -- (clip.west);
  \draw[cfarr] (clip.east) -- (pkt.west);

  % the three states feed the roll-up: level exits onto a right-margin bus,
  % one drop into a clear channel between the grid and the roll-up row, and a
  % single arrow up into the roll-up. No feed line crosses a box.
  \def\xret{13.45}\def\yret{-8.30}
  \foreach \b/\y in {b1/\rya, b2/\ryb, b3/\ryc}{
    \draw[cfline] (\b.east) -- (\xret,\y);
    \node[cfdot] at (\xret,\y) {};
  }
  \draw[cfline] (\xret,\rya) -- (\xret,\ryc);      % vertical return bus
  \coordinate (chan) at (\xret,\yret);             % channel, right end
  \draw[cfline] (\xret,\ryc) -- (chan) -- (roll.north |- chan);
  \draw[cfarr]  (roll.north |- chan) -- (roll.north);

\end{tikzpicture}
\endgroup

  \caption{Overview of \PhyReAct{} relative to question decomposition. The
  numbered views magnify rotation, floor-contact deformation, generated wall
  penetration, and the counterfactual rigid-wall rebound; the waveform shares
  the visual event clock. The scene is illustrative rather than a recorded
  execution.}
  \label{fig:teaser}
\end{figure}
\FloatBarrier

Despite this progress, recent evaluators capture complementary components of
physical diagnosis rather than a single auditable evidence chain. They provide
rule-level judgments~\cite{videophy2_2025,lin2026phyground}, temporal
localization~\cite{physioneval2026,chinchure2026spotlight}, measurement-grounded
tests in controlled settings~\cite{wang2026gauge}, and audio-physical or
cross-modal consistency checks~\cite{xie2025phyavbench,zhou2026avgenbench}.
These capabilities are nevertheless typically organized as benchmark-specific
protocols built around curated datasets, predefined phenomena, or manually
authored rubrics. The recent AV-Phys Agent is a particularly close point of
comparison: it combines a reasoning--action loop with deterministic acoustic
tools, but its released protocol remains tied to prompt-specific human-authored
rubrics and multimodal-model binary judgments~\cite{cui2026avphys}. The
ball--wall schematic in \cref{fig:teaser}, which illustrates the interface
rather than a recorded execution, makes the resulting integration gap concrete:
a defensible verdict must establish that the ball and wall are present,
localize floor- and wall-contact windows, measure trajectory and relative
geometry, and determine whether any detected impact sound aligns with the
wall-contact window. Although audio events can be detected in parallel,
synchrony cannot be resolved until that visual window is available. Unusable
evidence leaves the dependent claim unknown; if no claim is contradicted, any
unknown makes the clip abstain, while a terminal infrastructure failure
remains a separate outcome. The gap targeted here is therefore not the absence
of any one component, but their integration within a single evaluator that
compiles the prompt into typed physical obligations and a statically validated
plan, lets observations gate and scope already-declared specialist calls, and
uses fixed rules to compose usable evidence records---typed measurements or
explicitly tagged learned states---into three-valued decisions with traceable
time windows, tool outputs, values, rules, and provenance.

\PhyReAct{} addresses these requirements as a single physical-verification
system coupling a text-only planner, a video-aware semantic verifier, frozen
specialist operators, and fixed composition. For open-ended prompts that provide
no evaluation rubric, the planner operates before any frame is seen, compiling
exact prompt spans into typed physical obligations and a statically validated
execution plan. For evidence distributed across entities, events, time, and
modalities, the verifier reads timestamped frames densely on a shared A/V clock;
its observations activate, skip, or localize only calls already declared by the
plan, and those scoped calls dispatch SAM~3 grounding and
tracking~\cite{sam3_2025}, identity-based counting, eleven track-based physical
measurements, monocular depth, optical character recognition, and audio-event
detection~\cite{flexsed2025}. For heterogeneous or unusable outputs, each action
emits a claim-bound record containing either a typed measurement or an explicitly
tagged learned base state, together with its realized scope,
measured/abstained/errored status, artifact hashes, cost, and provenance. Typed
resolvers and fixed composition map usable records to supported, contradicted,
or unknown claim states and then to plausible, implausible, or abstain at clip
level; unavailable evidence remains unknown, while terminal infrastructure
failures remain separate. For opaque verdicts, the system returns an auditable
claim-level trace linking each decision to its prompt span, planned calls,
localized evidence, and composition rule. Together, these components turn an
open prompt and generated A/V clip into a dependency-aware evidence chain rather
than an unordered set of model answers. \Cref{fig:teaser} contrasts this design
with the evaluated question-decomposition baseline, which asks one closed
yes/no question per claim without a measurement trace.

To evaluate the critic at the granularity of its trace, we assemble a corpus of
$1{,}500$ generated clips with $2{,}582$ human-written, prompt-grounded flaw
records. Each record quotes the violated prompt span and provides a rationale,
severity, confidence, and, when available, a temporal span and object track;
these records are the only ground truth used here. A flaw-level protocol scores
each system finding as whole, partial, or missed against the corresponding
human record. On a $149$-clip recall analysis, \PhyReAct{} accounts for $228$ of
$304$ annotated defects, compared with $164$ for a published
question-decomposition evaluator given the same clips, claims, and served model.
A monolithic prompt to that same backbone accounts for $222$, so recall alone
does not isolate the value of the agentic organization; the principal
distinction is that every \PhyReAct{} decision retains its supporting evidence
and provenance. This analysis is recall-only, single-annotator, and not held
out, and therefore characterizes the current system rather than its
generalization (\cref{sec:eval}).

\PhyReAct{} casts physically grounded generation of, and a verdict on, a clip $\mathcal{V}$ for its
prompt $p$ as two maps we build and benchmark separately: a \emph{generation} map
$\hat{\mathcal{V}}\sim\operatorname{Gen}(p,C)$, the control-conditioned video diffusion of
\cref{sec:generation}, and an \emph{evaluation} map $Y=\operatorname{Crit}(\mathcal{V},p)$, the
auditable critic of \cref{sec:critic}. Both maps and the clip-level verdict are made precise there.

In summary, our contributions are threefold: (i) \PhyReAct{}, an integrated
physical-verification system that compiles prompts into typed obligations,
scopes specialist operators with observations, and returns three-valued
decisions with localized evidence and provenance; (ii) the $1{,}500$-clip corpus
and its flaw-level matching protocol; and (iii) a simulation-conditioned
generation testbed in which known MuJoCo geometry is rendered as depth or
silhouette control for a frozen Wan~2.2-VACE generator
~\cite{todorov2012mujoco,wan2025,vace2025}. The current report evaluates the
critic and generator separately. Together, these contributions provide an
auditable evaluation stack and a concrete interface for future physical
refinement.
